% Part: lambda-calculus
% Chapter: introduction
% Section: computable-lambda

\documentclass[../../../include/open-logic-section]{subfiles}

\begin{document}

\olfileid[fr]{lam}{int}{lrp}
\olsection{Les fonctions calculables sont \texorpdfstring{$\lambd$}{λ}-définissables}

\begin{thm}
\ollabel{thm:computable-lambda}
Toute fonction partielle calculable est
$\lambd$-définissable.
\end{thm}

\begin{proof}
Par l'équivalence des modèles de calcul, il suffit de traiter les
fonctions récursives partielles. Le
\olref[lam][rep][par]{thm:partial-recursive-lambda} les représente
par des lambda-termes au moyen de la forme normale de Kleene.
Les sections qui suivent donnent les constructions élémentaires
utilisées pour les fonctions récursives primitives.
\end{proof}

Nous emploierons désormais une notation plus
classique pour rendre la suite de la preuve plus
lisible. Nous écrirons, par exemple, $M(x,y,z)$
au lieu de $Mxyz$. Cette notation est suggestive,
mais rappelons que les termes du lambda-calcul non
typé n'ont pas d'arité attachée : pour le même terme~$M$,
$M(x,y)$ et $M(x,y,z,w)$ ont tout autant de sens.
Écrire $M(x,y,z)$ indique toutefois que nous considérons
$M$ comme une fonction de trois arguments et rend plus
claire l'intention des définitions. De même, nous dirons
« définissons~$M$ par $M(x,y,z)=\dots$ » plutôt que
« définissons~$M$ par
$M=\lambd[x][\lambd[y][\lambd[z][\dots]]]$ ».

\end{document}
